\paper{数学 B 类}

\begin{problems}

\item[一、]（本题 15 分）设球面 $S: x^2+y^2+z^2=1$，求以点 $M_0(0,0,a)$（$a\in\R$，$|a|>1$）为顶点的与 $S$ 相切的锥面方程．

\item[二、]（本题 15 分）设 $B\subset\R^n\ (n\ge2)$ 是单位开球，函数 $u,v$ 在 $\bar B$ 上连续，在 $B$ 内二阶连续可导，满足
\[
\begin{gathered}
  -\Delta u-(1-u^2-v^2)u=0,\qquad x\in B,\\
  -\Delta v-(1-u^2-v^2)v=0,\qquad x\in B,\\
  u(x)=v(x)=0,\qquad x\in\partial B,
\end{gathered}
\]
其中 $x=(x_1,x_2,\cdots,x_n)$，$\Delta u=\dfrac{\partial^2u}{\partial x_1^2}+\dfrac{\partial^2u}{\partial x_2^2}+\cdots+\dfrac{\partial^2u}{\partial x_n^2}$，$\partial B$ 表示 $B$ 的边界．证明：$u^2(x)+v^2(x)\le1\ (x\in\bar B)$．

\item[三、]（本题 15 分）设 $f(x)=x^{2021}+a_{2020}x^{2020}+a_{2019}x^{2019}+\cdots+a_2x^2+a_1x+a_0$ 为整系数多项式，$a_0\ne0$．设对任意 $0\le k\le2020$ 有 $|a_k|\le40$，证明：$f(x)=0$ 的根不可能全为实数．

\item[四、]（本题 20 分）设 $R=\{0,1,-1\}$，$S$ 为 $R$ 上的 3 阶行列式全体，即 $S=\{\det(a_{ij})_{3\times3}\mid a_{ij}\in R\}$．证明：$S=\{-4,-3,-2,-1,0,1,2,3,4\}$．

\item[五、]（本题 15 分）设函数 $f$ 在 $[-1,1]$ 内有定义，在 $x=0$ 的某邻域内连续可导，且 $\lim\limits_{x\to0}\dfrac{f(x)}{x}=a>0$．证明：级数 $\sum\limits_{n=1}^{\infty}(-1)^nf\Bigl(\dfrac{1}{n}\Bigr)$ 收敛，$\sum\limits_{n=1}^{\infty}f\Bigl(\dfrac{1}{n}\Bigr)$ 发散．

\item[六、]（本题 20 分）设 $f(x)=\ln\sum\limits_{n=1}^{\infty}\dfrac{\ee^{nx}}{n^2}$．证明函数 $f$ 在 $(-\infty,0)$ 内为严格凸的，并且对任意 $\xi\in(-\infty,0)$，存在 $x_1,x_2\in(-\infty,0)$ 使得
\[
  f'(\xi)=\frac{f(x_2)-f(x_1)}{x_2-x_1}.
\]
（称 $(a,b)$ 内的函数 $f$ 为严格凸的，如果对任何 $\alpha\in(0,1)$，以及 $x,y\in(a,b)$，$x\ne y$ 成立 $f(\alpha x+(1-\alpha)y)<\alpha f(x)+(1-\alpha)f(y)$．）

\end{problems}
